\chapter{\ifenglish Background Knowledge and Theory\else หลักการและทฤษฎี\fi}

\hspace{1.5em}บทนี้นำเสนอหลักการ แนวคิด และทฤษฎีที่เกี่ยวข้องกับการออกแบบและพัฒนาเกมในโครงงาน \\โดยมุ่งเน้นการอธิบายพื้นฐานทางด้านวิศวกรรมคอมพิวเตอร์และการออกแบบเกม ที่ถูกนำมาใช้ในการพัฒนาระบบต่าง ๆ ภายในเกม

\section{Roguelike Game Design}
\hspace{1.5em}Roguelike เป็นแนวเกมที่มีต้นกำเนิดจากเกม Rogue โดยมีลักษณะเด่นคือการเล่นแบบวนลูป และการสร้างประสบการณ์ที่แตกต่างในแต่ละรอบการเล่นผ่านระบบสุ่ม และความก้าวหน้าของตัวละคร [2][6]\par
\vspace{1\baselineskip}
ลักษณะสำคัญของเกมแนว Roguelike ได้แก่
\begin{itemize}
    \item \textbf{Permadeath}\newline
    \hspace{1.5em}เมื่อผู้เล่นแพ้หรือเสียชีวิต จะต้องเริ่มต้นใหม่ตั้งแต่ต้น ซึ่งช่วยเพิ่มความท้าทายและความตึงเครียดในการเล่น
    \item \textbf{Procedural Generation}\newline
    \hspace{1.5em}การสร้างฉาก ศัตรู หรือเหตุการณ์แบบสุ่ม ทำให้แต่ละรอบการเล่นมีความแตกต่างกัน
    \item \textbf{Core Gameplay Loop ที่ชัดเจน}\newline
    \hspace{1.5em}ผู้เล่นจะทำกิจกรรมซ้ำในรูปแบบลูป เช่น การสำรวจ ต่อสู้ เก็บทรัพยากร และพัฒนาตัวละคร
    \item \textbf{Progression System}\newline
    \hspace{1.5em}ผู้เล่นสามารถพัฒนาความสามารถของตัวละครผ่านการเล่นซ้ำ เช่น การปลดล็อกสกิลหรือไอเทม
\end{itemize}

ในโครงงานนี้ แนวคิด Roguelike ถูกนำมาประยุกต์ใช้ร่วมกับระบบการทำอาหาร โดยผู้เล่นจะต้องออกล่ามอนสเตอร์เพื่อเก็บวัตถุดิบ นำมาปรุงอาหาร และใช้ผลลัพธ์จากอาหารในการพัฒนาความสามารถของ\\ตัวละคร ซึ่งสร้างเป็น Core Gameplay Loop ดังนี้

\begin{center}
  \textbf{ล่า $\rightarrow$ ได้วัตถุดิบ $\rightarrow$ ทำอาหาร $\rightarrow$ แข่งขัน $\rightarrow$ พัฒนา $\rightarrow$ วนซ้ำ}
\end{center}

แนวทางนี้ช่วยให้ระบบการทำอาหารถูกผนวกเข้ากับกลไกหลักของเกมอย่างเป็นธรรมชาติ \\และทำให้ผู้เล่นเกิดการเรียนรู้ผ่านการเล่นซ้ำ ซึ่งสอดคล้องกับวัตถุประสงค์ของโครงงาน

\newpage
\section{Game Loop Design}
\hspace{1.5em}Game Loop เป็นแนวคิดพื้นฐานในการออกแบบเกม โดยหมายถึงลำดับของกิจกรรมหลักที่ผู้เล่นจะทำซ้ำอย่างต่อเนื่องตลอดการเล่นเกม (core repetitive cycle) การออกแบบ Game Loop ที่ชัดเจนและมีความสมดุลจะช่วยให้เกมมีความน่าสนใจ และสามารถรักษาการมีส่วนร่วมของผู้เล่นได้ (player engagement)

โดยทั่วไป Game Loop จะประกอบด้วยขั้นตอนหลัก เช่น การสำรวจ (exploration) การต่อสู้ \\(combat) การเก็บทรัพยากร (resource collection) และการพัฒนาตัวละคร (progression) ซึ่งจะถูกออกแบบให้เชื่อมโยงกันเป็นวงจร [1][5]

ในโครงงานนี้ ได้มีการออกแบบ Core Gameplay Loop โดยผสมผสานระบบการต่อสู้และการทำอาหารเข้าด้วยกัน ดังนี้

\begin{center}
\textbf{ล่ามอนสเตอร์ $\;\rightarrow\;$ ได้วัตถุดิบ $\;\rightarrow\;$ ทำอาหาร $\;\rightarrow\;$ แข่งขัน $\;\rightarrow\;$ พัฒนาตัวละคร $\;\rightarrow\;$ วนซ้ำ}
\end{center}

การออกแบบลูปดังกล่าวทำให้ผู้เล่นมีเป้าหมายที่ชัดเจนในแต่ละรอบของการเล่น โดยการล่ามอนสเตอร์เพื่อเก็บวัตถุดิบจะเชื่อมโยงไปสู่ระบบการทำอาหาร ซึ่งผลลัพธ์จากการทำอาหารจะส่งผลต่อความสามารถ\\ของตัวละคร และนำไปสู่การแข่งขันเพื่อปลดล็อกเนื้อหาใหม่

นอกจากนี้ เกมยังมีการออกแบบลูปย่อย (Secondary Loops) เช่น
\begin{itemize}
  \item การพัฒนาความชำนาญของเมนู (Food Mastery Loop)
  \item การปลดล็อกสกิลจากอาหาร (Skill Progression Loop)
  \item การดำเนินเนื้อเรื่อง (Story Progression Loop)
\end{itemize}

ลูปย่อยเหล่านี้ช่วยเพิ่มความหลากหลายของการเล่น และสนับสนุนให้ผู้เล่นกลับมาเล่นซ้ำ (replayability) ซึ่งเป็นองค์ประกอบสำคัญของเกมแนว Roguelike [2]

โดยสรุป การออกแบบ Game Loop ในโครงงานนี้ช่วยเชื่อมโยงระบบต่าง ๆ ภายในเกมเข้าด้วยกันอย่างเป็นระบบ และสร้างประสบการณ์การเล่นที่ต่อเนื่องและมีความหมายต่อผู้เล่น

\section{Behaviors in Games}
\hspace{1.5em}พฤติกรรมของตัวละคร เป็นองค์ประกอบสำคัญในการออกแบบเกม โดยเฉพาะในระบบที่มีตัวละครที่\\ไม่ใช่ผู้เล่น (Non-Player Characters: NPCs) ซึ่งจำเป็นต้องมีการตัดสินใจและตอบสนองต่อสถานการณ์ต่าง ๆ ภายในเกมอย่างเหมาะสม [3]

การออกแบบ Behavior ที่ดีจะช่วยให้ตัวละครมีความสมจริง และเพิ่มความท้าทายให้กับผู้เล่น

ในโครงงานนี้ ได้มีการเลือกใช้รูปแบบของ Behavior ที่แตกต่างกันตามบทบาทของตัวละคร \\เพื่อให้เหมาะสมกับลักษณะการทำงานของแต่ละประเภท โดยสามารถแบ่งออกเป็น 3 รูปแบบหลัก ดังนี้ [4]
\newpage
\subsection{Model-Based Reflex Agent}
\hspace{1.5em}Model-Based Reflex Agent เป็นรูปแบบที่ใช้สถานะภายใน (internal state) เพื่อช่วยในการ\\ตัดสินใจ โดยตัวละครจะสามารถจดจำข้อมูลบางอย่างของโลก เช่น ตำแหน่งล่าสุดของผู้เล่น

ในเกมนี้ รูปแบบดังกล่าวถูกนำมาใช้กับมอนสเตอร์ โดยมอนสเตอร์สามารถติดตามผู้เล่น โจมตี และค้นหาผู้เล่นเมื่อหลุดจากการมองเห็น ซึ่งช่วยให้พฤติกรรมมีความต่อเนื่องและสมจริงมากขึ้น

\subsection{Goal-Based Agent}
\hspace{1.5em}Goal-Based Agent เป็นรูปแบบที่มีเป้าหมาย เป็นตัวกำหนดการตัดสินใจ โดยระบบจะเลือกการ\\กระทำที่ช่วยให้บรรลุเป้าหมายได้ดีที่สุด

ในโครงงานนี้ แนวคิดดังกล่าวถูกนำมาใช้กับตัวละครคู่แข่งในการแข่งขันทำอาหาร ซึ่งมีเป้าหมายหลักคือการชนะการแข่งขัน โดยจะเลือกการกระทำ เช่น การเก็บวัตถุดิบ หรือการเร่งทำอาหาร เพื่อเพิ่มโอกาสในการชนะ

\subsection{Utility-Based Agent}
\hspace{1.5em}Utility-Based Agent เป็นรูปแบบ ที่ใช้การคำนวณค่าความเหมาะสม (utility) ของแต่ละการกระทำ และเลือกการกระทำที่มีค่ามากที่สุด

ในเกมนี้ ใช้กับตัวละครผู้ช่วย (Companion) โดยระบบจะพิจารณาสถานการณ์ เช่น ระดับพลังชีวิตของผู้เล่น หรือความอันตรายของศัตรู เพื่อเลือกการกระทำที่เหมาะสม เช่น การรักษาผู้เล่น หรือการเบี่ยงเบนความสนใจของมอนสเตอร์

โดยสรุป การเลือกใช้ Behavior ที่แตกต่างกันตามบทบาทของตัวละครช่วยให้ระบบมีความยืดหยุ่น และสามารถสร้างประสบการณ์การเล่นที่หลากหลายและสมจริงมากยิ่งขึ้น

\section{Pathfinding Algorithm}
\hspace{1.5em}Pathfinding เป็นกระบวนการคำนวณเส้นทางที่เหมาะสมที่สุดจากจุดเริ่มต้นไปยังจุดหมายปลายทาง \\โดยคำนึงถึงข้อจำกัดของสภาพแวดล้อม เช่น สิ่งกีดขวาง หรือพื้นที่ที่ไม่สามารถเคลื่อนที่ผ่านได้ ซึ่งเป็นองค์ประกอบสำคัญในการพัฒนาพฤติกรรมของตัวละครในเกม [3]

ในโครงงานนี้ ได้เลือกใช้ A* (A-star) Algorithm ซึ่งเป็นอัลกอริทึมที่ได้รับความนิยมในการค้นหาเส้นทาง เนื่องจากสามารถหาทางที่มีประสิทธิภาพและมีความแม่นยำสูง โดย A* จะพิจารณาค่าใช้จ่ายรวม \\(cost function) ดังนี้ [10][11]

\[
f(n) = g(n) + h(n)
\]

โดยที่
\begin{itemize}
  \item $g(n)$ คือ ค่าใช้จ่ายจากจุดเริ่มต้นไปยังโหนดปัจจุบัน
  \item $h(n)$ คือ ค่า heuristic ซึ่งเป็นค่าประมาณระยะทางจากโหนดปัจจุบันไปยังจุดหมาย
\end{itemize}

การใช้ heuristic ช่วยให้ A* สามารถค้นหาเส้นทางได้รวดเร็วกว่าวิธีการแบบ brute-force และลด\\จำนวนโหนดที่ต้องพิจารณา

ในการพัฒนาเกมนี้ A* ถูกนำมาใช้ร่วมกับระบบ NavMesh ของ Unity ซึ่งทำหน้าที่สร้างโครงสร้างพื้นที่ที่ตัวละครสามารถเคลื่อนที่ได้ (navigation mesh) และช่วยให้ตัวละครสามารถหลีกเลี่ยงสิ่งกีดขวางได้อย่างอัตโนมัติ

กระบวนการทำงานโดยรวมคือ ตัวละครจะกำหนดตำแหน่งเป้าหมาย จากนั้นระบบ NavMesh จะใช้ A* Algorithm ในการคำนวณเส้นทางที่เหมาะสม และส่งผลลัพธ์ไปยังระบบการเคลื่อนที่ของตัวละคร

การเลือกใช้ A* Algorithm ร่วมกับ NavMesh ช่วยให้การเคลื่อนที่ของมอนสเตอร์มีความสมจริง และสามารถตอบสนองต่อสภาพแวดล้อมภายในเกมได้อย่างมีประสิทธิภาพ

\section{Mathematical Model for Scoring}
\hspace{1.5em}ในการออกแบบระบบการแข่งขันทำอาหาร โครงงานนี้มุ่งเน้นการประเมินคุณภาพของอาหารในเชิง \\“ความใกล้เคียงของรสชาติ” ระหว่างอาหารที่ผู้เล่นสร้างและอาหารเป้าหมาย โดยใช้แบบจำลองทาง\\คณิตศาสตร์เพื่อให้การให้คะแนนมีความเป็นระบบและตรวจสอบได้

\subsection{การแทนค่ารสชาติด้วยเวกเตอร์ (Flavor Vector)}
\hspace{1.5em}อาหารแต่ละจานถูกแทนในรูปของเวกเตอร์ 5 มิติ ได้แก่

\begin{center}
\textbf{Flavor Vector} = (spicy, salty, sweet, sour, bitter)
\end{center}

ซึ่งสะท้อนถึงองค์ประกอบของรสชาติในอาหาร

\subsection{Manhattan Distance}
\hspace{1.5em}ในการวัดความแตกต่างของรสชาติ โครงงานนี้เลือกใช้ Manhattan Distance ซึ่งมีสมการดังนี้ [12][13]

\[
\text{Distance} = \left|P_1 - T_1\right| + \left|P_2 - T_2\right| + \left|P_3 - T_3\right| + \left|P_4 - T_4\right| + \left|P_5 - T_5\right|
\]

โดยที่ $P$ คือเวกเตอร์ของอาหารที่ผู้เล่นสร้าง และ T คือเวกเตอร์ของอาหารเป้าหมาย

\vspace{1\baselineskip}
วิธีนี้สามารถวัด “ความแตกต่างเชิงปริมาณ” ของแต่ละรสชาติได้โดยตรง และสะท้อนความผิดพลาดในแต่ละมิติของรสชาติอย่างชัดเจน
\newpage
\subsection{เหตุผลในการเลือก Manhattan Distance}
\hspace{1.5em}การเลือกใช้ Manhattan Distance แทนวิธีอื่น เช่น Cosine Similarity หรือ Euclidean Distance มีเหตุผลดังนี้

\begin{itemize}
  \item \textbf{สะท้อนความแตกต่างของรสชาติแบบตรงไปตรงมา}\newline
  \hspace{1.5em}Manhattan Distance คำนวณจากผลรวมของความต่างในแต่ละมิติ ทำให้สามารถตีความได้ง่าย เช่น รสใดที่มากหรือน้อยเกินไป
  \item \textbf{เหมาะกับระบบที่ต้องการความแม่นยำของค่าจริง (absolute value)}\newline
  \hspace{1.5em}ในเกมนี้ ผู้เล่นต้องทำอาหารให้ “ใกล้ค่าที่กำหนด” ไม่ใช่แค่มีสัดส่วนที่คล้ายกัน
  \item \textbf{มีความซับซ้อนในการคำนวณต่ำ (computational efficiency)}\newline
  \hspace{1.5em}เนื่องจากไม่มีการยกกำลังหรือหาร ทำให้เหมาะกับการใช้งานในระบบที่ต้องคำนวณแบบ real-time
\end{itemize}

\subsection{เปรียบเทียบกับ Cosine Similarity}
\hspace{1.5em}Cosine Similarity เป็นวิธีที่ใช้วัดความคล้ายคลึงของเวกเตอร์ในเชิงทิศทาง (direction) โดยไม่สนใจขนาดของเวกเตอร์ [14]

\vspace{1\baselineskip}
อย่างไรก็ตาม วิธีนี้ไม่เหมาะกับระบบของเกม เนื่องจาก:

\begin{itemize}
  \item อาหารที่มี “สัดส่วนรสชาติคล้ายกัน” แต่มีความเข้มต่างกัน จะได้คะแนนสูง \\เช่น (5,3,2) และ (10,6,4) จะถือว่าเหมือนกัน
  \item ไม่สามารถสะท้อนความแตกต่างเชิงปริมาณของรสชาติได้อย่างชัดเจน
\end{itemize}

\vspace{0.2\baselineskip}
ดังนั้น Cosine Similarity จึงไม่เหมาะกับการใช้งานในกรณีที่ต้องการวัด “ความใกล้เคียงของค่าจริง”

\subsection{การคำนวณคะแนน}
\hspace{1.5em}คะแนนรสชาติถูกคำนวณจากสมการ

\[
\text{Flavor Score} = \text{MaxScore} - k \cdot \text{Distance}
\]

และนำไปรวมกับปัจจัยอื่นเพื่อให้ได้คะแนนสุดท้าย

\[
\text{Final Score} = \text{Flavor Score} + \text{Cooking Skill Score} + \text{Ingredient Bonus}
\]

\subsection{สรุป}
\hspace{1.5em}Manhattan Distance เป็นวิธีที่เหมาะสมสำหรับระบบนี้ เนื่องจากสามารถสะท้อนความแตกต่างของรสชาติได้อย่างชัดเจน มีความซับซ้อนต่ำ และสอดคล้องกับรูปแบบ gameplay ที่ต้องการให้ผู้เล่นปรับค่า\\รสชาติให้ใกล้เคียงกับเป้าหมายมากที่สุด
